\documentclass[10pt,a4paper]{amsart}
\usepackage{amsmath,amssymb,amsthm,graphicx}
\usepackage{amsmath}
\usepackage{amssymb}
\usepackage{hyperref}
\newtheorem{Theorem}{Theorem}[section]
\newtheorem{Corollary}[Theorem]{Corollary}
\newtheorem{Lemma}[Theorem]{Lemma}
\newtheorem{Proposition}[Theorem]{Proposition}
{ \theoremstyle{definition}
	\newtheorem{Definition}[Theorem]{Definition}
	\newtheorem{Note}[Theorem]{Note}
	\newtheorem{Example}[Theorem]{Example}
	\newtheorem{Remark}[Theorem]{Remark} }

\usepackage[all]{xy}


\def\KK{\ensuremath\mathbb{K}}
\def\MM{\ensuremath\mathbb{M}}
\def\BB{\ensuremath\mathbb{B}}
\def\GL{\ensuremath\mathbf{GL}}
\def\UU{\ensuremath\mathbf{U}}
\def\PP{\ensuremath\mathbb{P}}
\def\CC{\ensuremath{\mathbb{C}}}
\def\NN{\ensuremath{\mathbb{N}}}
\def\QQ{\ensuremath{\mathbb{Q}}}
\def\RR{\ensuremath{\mathbb{R}}}
\def\ZZ{\ensuremath{\mathbb{Z}}}

\newcommand{\profinQ}{\mathfrak{Q}}
\newcommand{\reals}{\mathbb{R}}
\newcommand{\complexs}{\mathbb{C}}
\newcommand{\naturals}{\mathbb{N}}
\newcommand{\integers}{\mathbb{Z}}
\newcommand{\rationals}{\mathbb{Q}}
\newcommand{\quaternions}{\mathbb{H}}
\newcommand{\hyperbolic}{\mathbb{H}}
\newcommand{\K}{\mathbb{K}}
\newcommand{\realsorcomplexs}{\mathbb{K}}
\newcommand{\finitefield}{\mathbb{F}}
\newcommand{\Proj}{\mathbb{P}}
\DeclareMathOperator{\id}{id}
\DeclareMathOperator{\Ind}{Ind}
\newcommand{\boundary}[1]{\partial#1}
\newcommand{\boundedops}{\mathcal{B}}
\newcommand{\bundleover}{\!\!\downarrow\!\!}
\newcommand{\abs}[1]{\left\lvert#1\right\rvert} %absolute value
%\newcommand{\bigabs}[1]{\left\lvert#1\right\rvert} %absolute value
\newcommand{\norm}[1]{\left\lVert#1\right\rVert}
\newcommand{\generate}[1]{\langle#1\rangle}
\newcommand{\tensor}{\otimes}
\newcommand{\into}{\hookrightarrow}
\newcommand{\onto}{\twoheadrightarrow}
\newcommand{\iso}{\cong}
\newcommand{\disjointunion}{\amalg}
\newcommand{\subgroup}{\leq}
\newcommand{\supergroup}{\geq}
\newcommand{\superset}{\supset}
\newcommand{\normalsubgroup}{\lhd}
\newcommand{\semiProd}{\rtimes}
\newcommand{\semiprod}{\semiProd}
\newcommand{\dividesnot}{\nmid}
\newcommand{\divides}{\mid}
\DeclareMathOperator{\Pos}{Pos}
\DeclareMathOperator{\Aut}{Aut}
\DeclareMathOperator{\supp}{supp} %support
\DeclareMathOperator{\clos}{clos} %closure
\DeclareMathOperator{\im}{im} %image
\DeclareMathOperator{\vol}{vol} %Volume
\DeclareMathOperator{\diam}{diam} %Diameter
\DeclareMathOperator{\dist}{dist} %Distance
\DeclareMathOperator{\ord}{ord} %order
\DeclareMathOperator{\End}{End} %Endomorphisms
\DeclareMathOperator{\Hom}{Hom} %Homomorphisms
\DeclareMathOperator{\Span}{span}
\DeclareMathOperator{\spin}{spin}
\DeclareMathOperator{\diag}{diag}
\DeclareMathOperator{\spec}{Spec}
\DeclareMathOperator{\trace}{trace}
\DeclareMathOperator{\tr}{tr}
\DeclareMathOperator{\pr}{pr}
%\DeclareMathOperator{\res}{res} %Restriction
\DeclareMathOperator{\grad}{grad} %Restriction
% \DeclareMathOperator{\Sp}{Sp} %Spur % schon definiert???
\DeclareMathOperator{\fl}{f\/l} %flip
\DeclareMathOperator{\coker}{coker}
\DeclareMathOperator{\Char}{Char}
\DeclareMathOperator{\ind}{ind}
\DeclareMathOperator{\sgn}{sgn}
\DeclareMathOperator{\Ric}{Ric} % Ricci curvature
\DeclareMathOperator{\scal}{scal} % Scalar curvature
\DeclareMathOperator*{\dirlim}{dirlim}
\DeclareMathOperator*{\invlim}{invlim}
\DeclareMathOperator{\rank}{rank}
\DeclareMathOperator{\coind}{coind}
\DeclareMathOperator{\res}{res}
\newcommand{\bigdisjointunion}{\bigsqcup}
\DeclareMathOperator{\Orb}{Orb}
\DeclareMathOperator{\GTOP}{G-TOP}
\newcommand{\forget}[1]{}
\newenvironment{Forget}{}{}
\newcommand{\innerprod}[1]{\langle #1 \rangle}
\title{A morphism of groupoids bijective on orbits and on isotropy groups is an equivalence}
\author{Massimiliano Puglisi, Thomas Schick, Vito Felice Zenobi}
\date{doi:10.3842/SIGMA.2025.093; CC BY 4.0}
\begin{document}
\maketitle
\noindent\textit{Editorial note.} All delivered mathematical text is the author's own, copied from the pinned publisher source sources/geometry-fl04-stolz-positive-scalar-curvature-st/pinned-original-source.tex; no statement or proof step was rewritten, reordered or removed.\par\medskip
\noindent\textit{Editorial note.} Preamble provenance: the retained preamble is the source's own layout adapter, if any, plus the author's macro block at lines 11-108 of sources/geometry-fl04-stolz-positive-scalar-curvature-st/pinned-original-source.tex (everything between the publisher class metadata and \textbackslash{}begin\{document\}, comments excluded). Nothing was added, no macro was invented, and the publisher document class is deliberately not reproduced; the excerpt serves only as preamble, never as a span coordinate.\par\medskip
\noindent\textit{Editorial note.} Metadata provenance: source\_id and source\_version were taken from sources/geometry-fl04-stolz-positive-scalar-curvature-st/license-observation.json; the title from SOURCE.json; the author list from the source metadata record; the source URL from the first non-empty of source\_url, source\_urls[0] and the evidence's fixed\_version\_url/license\_url. The licence evidence file used by this card is sources/geometry-fl04-stolz-positive-scalar-curvature-st/license-observation.json, recorded in license.evidence\_path so that the binding target is explicit.\par\medskip
\noindent\textit{Editorial note.} External references are navigation pointers only; every deliberately kept step of the author's argument stays inside the counted spans. The displayed numbers of those pointers are the author's own numbering under the author's \textbackslash{}newtheorem declarations (section.index as declared), computed mechanically by the lane's scanner.\par\medskip
\noindent\textit{Editorial note.} The counted proof contains no reference to any labelled item and no bibliography citation, so the card needs neither a navigation stub nor a bibliography excerpt.\par\medskip
\noindent\textit{Editorial note.} The statement is labelled by the author before the environment body, which is why the label sits at source line 1261.\par\medskip
\noindent\textit{Allowed standard prerequisites (editorial).}\begin{itemize}
\item The orbit set pi\_0 of a groupoid and the isotropy groups of its objects. Recalled inside the counted proof at lines 1270-1276. Fixes the two invariants whose bijectivity the hypothesis asserts.
\item A disjoint union of equivalences of groupoids is again an equivalence. Used inside the counted proof at lines 1276-1278. Turns the displayed decomposition of the groupoid into isotropy groups into the claimed equivalence.
\end{itemize}\medskip
	\begin{Lemma}\label{lem:groupoid_eq}
		Let $T\colon \mathcal{G}\to \mathcal{H}$ be a morphism of groupoids such that it induces a
		bijection between the sets of orbits and for each object $x\in \mathcal{G}$ an
		isomorphism of isotropy groups \smash{$\mathcal{G}_x^x\to \mathcal{H}_{T(x)}^{T(x)}$}. Then~$T$
		is an equivalence of groupoids.
	\end{Lemma}

	\begin{proof}
		The assertion follows directly from the following commutative diagram:
		\begin{equation*}
			\xymatrix{\displaystyle\bigsqcup_{[x]\in\pi_0(\mathcal{G})}\mathcal{G}^x_x\ar[r]^{\cong}\ar[d]^(.6){\simeq} &
				\displaystyle\bigsqcup_{[x]\in\pi_0(\mathcal{H})}\mathcal{H}^{T(x)}_{T(x)}\ar[d]^(.6){\simeq}\\
				\mathcal{G}\ar[r]^T&\mathcal{H}.}
		\end{equation*}
		
		Here, we denote $\pi_0(\mathcal{G})$ the set of orbits of $\mathcal{G}$ and we use that $T$
		induces a bijection between~$\pi_0(\mathcal{G})$ and $\pi_0(\mathcal{H})$
		and that a disjoint union of equivalences of groupoids is again an equivalence
		of groupoids.
	\end{proof}

\end{document}
